\chapter{Produto Educacional e Sequência Didática}

\section{Concepção do produto educacional}

O produto educacional une uma plataforma interativa em português brasileiro a uma sequência didática para o Ensino Médio. A plataforma roda direto no navegador com cenários sintéticos, executando todos os cálculos no próprio computador. Seu acesso dispensa criação de conta ou cadastro de usuário, operando sem envio de respostas ou coleta de dados pessoais. Isso facilita o uso mesmo sob conectividade limitada, bastando abrir o recurso nos computadores ou no projetor. A ferramenta apoia a mediação docente em Computação, preservando o espaço da explicação e do diálogo conduzidos pelo professor.

A proposta tem cinco módulos introdutórios de \textit{machine learning}: regressão linear, regressão logística, \textit{k}-vizinhos mais próximos, árvore de decisão e \textit{k}-means. Longe de formarem um catálogo de comandos ou conteúdos soltos, eles apresentam tarefas distintas de previsão numérica e classificação, além da formação de grupos para que a turma compare o que muda entre eles. Todos ocupam lugar central na sequência, mobilizando camadas matemáticas próprias. Na regressão linear, o estudo da função e da linha de tendência fundamenta a previsão e o cálculo de resíduos ou erros, enquanto a logística articula probabilidade ao limiar de decisão. Já o \textit{k}-NN destaca a distância entre pontos, o \textit{k}-means recalcula centroides pela média e a árvore encadeia condições.

A diferença de formalização não diminui nenhum módulo. Ela decorre do repertório de estudantes iniciantes e da necessidade de tornar visível, no gráfico ou na regra, uma relação que a turma consiga discutir durante a aula.

Mais do que estarem na mesma plataforma, os modelos se conectam por perguntas recorrentes nos módulos. A turma identifica a origem do dado disponível e a saída esperada antes de testar qual ajuste pode alterá-la. A resposta varia em cada caso. Enquanto a reta gera uma previsão numérica, a regressão logística calcula uma probabilidade confrontada com um limiar. Por sua vez, o \textit{k}-NN consulta exemplos próximos e a árvore percorre regras condicionais, ao passo que o \textit{k}-means reorganiza pontos mesmo sem categorias prévias.

Com as mesmas perguntas em cada módulo, os estudantes comparam os modelos sem reduzi-los a uma lista de nomes. O professor percebe se a dificuldade está em um conceito específico ou na leitura de dados e gráficos, que atravessa todos os módulos.

\section{Organização pedagógica da plataforma}

Cada módulo parte de uma situação-problema e da origem dos dados para questionar o que o modelo tenta aprender. Antes de usar os controles, o estudante prevê uma consequência. Em seguida, modifica dados ou parâmetros definidos para o experimento, observando o efeito gerado no gráfico e na saída do modelo para explicar o que ocorreu. Esse movimento da previsão ao experimento, seguido pela observação e pela explicação, repete-se nos módulos para explicitar o raciocínio por trás de cada manipulação.

Mover a inclinação da reta leva o estudante a comparar a previsão com o valor observado, enquanto a troca de um limiar questiona por que a classe muda sem alterar a porcentagem calculada. A mesma lógica orienta o ajuste de vizinhos ou de divisões da árvore, assim como a posição dos centroides.

A sequência abre com recursos de inteligência artificial que os estudantes já usam e chega logo a um problema concreto com dados. A Matemática entra à medida que a leitura do gráfico e a decisão do modelo na situação exigem. Toda representação visual é acompanhada por tabelas e textos que explicitam a regra, evitando que a compreensão dependa apenas da imagem. As visualizações em 3D do protótipo entram como recurso complementar dispensável ao conceito central, mantendo na interface o uso por teclado e a opção de reduzir movimentos.

\section{Cenários e dados dos experimentos}

Os dados das atividades serão quase todos sintéticos. No segundo encontro, a turma poderá montar uma pequena tabela anônima com informações simples e voluntárias, ou mesmo medir objetos da sala. Essa tabela servirá apenas ao ensino, ficando fora do corpus. Os demais cenários descartam o uso de notas escolares, registros de comportamento ou qualquer dado pessoal dos estudantes.

Os cenários constituem uma proposta inicial sujeita a ajustes após a definição da escola e da turma. Para a regressão linear, a ideia é relacionar temperatura e consumo fictício de energia em uma sala, enquanto a regressão logística usa a leitura de um sensor simulado para discutir probabilidade e emissão de alerta. Já o \textit{k}-NN classificará situações descritas por duas medidas, ao passo que a árvore de decisão simulará o acionamento de uma irrigação conforme a umidade e a temperatura. No \textit{k}-means, pontos de coleta em um mapa fictício de bairro permitirão formar grupos não rotulados a partir da proximidade com os centroides. Assim, além do apelo visual, cada contexto deve tornar visível o efeito de alterar dados ou parâmetros em um nível compatível com a Matemática da série. Antes da aplicação, os exemplos serão revistos para adequar a linguagem e o tempo previsto ao perfil da turma.

\section{Sequência didática provisória}

A sequência foi planejada, por ora, para oito encontros de cerca de cinquenta minutos. No campo inicialmente previsto, o laboratório de informática da escola já conta com 20 computadores conectados à internet; caso a aplicação mude de unidade, a infraestrutura local precisará ser conferida. A confirmação da turma e do número de aulas dentro do componente curricular ainda depende da viabilidade do campo. O cenário preferencial é uma turma de Ensino Médio do próprio pesquisador, possivelmente na disciplina Projeto Permanente por Afinidade, podendo migrar para Matemática ou outro espaço curricular conforme a lotação e o calendário exigirem.

Todo encontro mantém a mesma estrutura, mesmo quando o modelo muda. A aula abre com uma situação ou pergunta que delimita o problema. Após registrar uma breve previsão individual, cada estudante explora os controles individualmente antes do debate coletivo e da síntese docente no fechamento. Essa rotina dá continuidade às aulas e integra a plataforma ao trabalho de construir e revisar hipóteses, evitando que o uso do site funcione como simples intervalo entre explicações.

O primeiro encontro serve para nivelar conceitos. Antes de qualquer modelo, a turma discute aplicações cotidianas de inteligência artificial. Depois distingue os dados e a regra do modelo, até apontar o que seria entregue ao sistema e o que se espera dele. O diagnóstico inicial é individual e acontece no começo desse encontro, antes da discussão, ou na aula anterior, conforme o tempo. Registra ideias de partida, sem antecipar os nomes técnicos que aparecem depois.

Do segundo ao sexto encontro entram, um por aula, os cinco módulos da plataforma, indo das regressões aos modelos de classificação e agrupamento. No segundo encontro, a organização da tabela da turma antecede a regressão linear. Cada aula retoma elementos anteriores para evitar que os modelos fiquem isolados em uma única experiência.

A regressão linear e a regressão logística abrem a parte dos modelos supervisionados. Na linear, a turma identifica a linha de tendência no plano cartesiano para fazer previsões e compará-las aos valores observados, introduzindo a noção de erro. A logística retoma o uso de exemplos conhecidos e direciona o debate para o cálculo de probabilidade e a escolha de classe pelo limiar. Essa sequência imediata permite contrastar a previsão de um valor numérico com a definição de uma categoria, mostrando que ambas dependem dos dados e das escolhas do modelo.

O \textit{k}-NN e a árvore de decisão são duas maneiras de chegar a uma classificação. No \textit{k}-NN, os estudantes observam o papel da distância e do número de vizinhos na votação; na árvore, acompanham regras condicionais e os caminhos até uma decisão. Comparar os dois afasta a ideia de que exista um algoritmo melhor que os outros. Ambos usam dados e regras, cada um à sua maneira, e entregam saídas que precisam ser lidas à luz do problema e dos exemplos disponíveis.

O \textit{k}-means traz uma mudança: os pontos não vêm com categorias dadas de antemão. A atividade volta à pergunta sobre o que se quer produzir com os dados e mostra que formar grupos por proximidade não é escolher uma classe conhecida. A atualização dos centroides e a observação da reorganização dos pontos poderão exigir mais tempo, pois a diferença entre agrupamento e classificação pode levar o professor a retomar a linguagem dos encontros anteriores. Uma síntese breve registra o modo como os cinco modelos articulam dados e regras para gerar uma saída, destacando o que muda conforme a finalidade de cada tarefa.

O sétimo encontro compara os cinco modelos. Diante de novas situações-problema, a turma justifica com base nos dados disponíveis se o caso exige previsão numérica, escolha de categoria ou formação de grupos. No oitavo encontro, cada estudante responde primeiro ao bloco de questões equivalentes às do diagnóstico e, na sequência, resolve um desafio de transferência inédito. Nessa etapa final, é preciso apontar os dados necessários para associar a tarefa a um modelo compatível, explicando o efeito esperado caso um parâmetro ou dado mude na regra.

O encontro comparativo recupera diferenças que possam ter ficado dispersas e prepara a transferência para o desafio final. Ao justificar que tipo de saída o problema demanda a partir dos dados disponíveis, o estudante considera a finalidade da tarefa antes de escolher o modelo. Enquanto o núcleo comparável permite examinar mudanças frente ao diagnóstico, o desafio de transferência mostra como os conceitos são mobilizados em novo contexto, entrando como evidência adicional sem equiparação direta com questões mais simples.

O trabalho na plataforma ocorrerá individualmente, com os computadores do laboratório e, conforme a disponibilidade confirmada no campo, Chromebooks para completar o acesso da turma. No segundo, no quarto e no sexto encontros, cada estudante também preencherá um registro curto antes e depois de explorar o site, cobrindo assim um exemplo de regressão, um de classificação e outro de agrupamento. Já o instrumento final será inteiramente individual. Caso o laboratório fique indisponível, o professor projetará os experimentos para a turma, preservando as pausas de previsão e explicação individual.

O roteiro preliminar dos oito encontros articula objetivos de aprendizagem e tempo previsto às ações de mediação com questões formativas, indicando ainda adaptações de infraestrutura e limites didáticos por módulo. Sem armazenamento de respostas digitadas no site, o corpus reunirá o diagnóstico inicial e o instrumento final completo, somados aos três registros individuais intermediários e às anotações docentes.

A versão do site usada no campo será conferida, antes da intervenção, com os recursos previstos na sequência. Se o cenário da plataforma tiver dados diferentes dos discutidos no roteiro da aula, a diferença será corrigida ou explicada aos estudantes; o mesmo vale para os textos de apoio e as representações matemáticas.

A pesquisa não pode analisar estudantes como se tivessem passado por uma sequência diferente do recurso usado. Essa conferência prática garante que a plataforma e os instrumentos avaliem exatamente o que foi trabalhado nas aulas.

\section{Materiais e limites do produto}

Além da plataforma, a aplicação contará com o diagnóstico inicial e três registros curtos de aula, somados às notas docentes e ao instrumento final com seu questionário de percepção. Nos registros intermediários, um campo recolhe a previsão prévia ao experimento e outro a explicação posterior. Já o questionário final consulta a turma sobre a clareza e o nível de dificuldade no manuseio do site, sem valer como evidência de aprendizagem.

O produto não comprova aprendizagem só por existir. Serve para que professor e estudantes discutam dados e parâmetros do modelo em sala.

A adequação didática do produto só poderá ser avaliada confrontando o planejamento com as condições reais do campo e as produções da turma, longe de constituir solução pronta para qualquer escola. Antes das aulas, um teste prático conferirá a abertura local do site nos computadores e a legibilidade da tela, checando ainda a resposta dos controles e a visibilidade no projetor. Caso algum equipamento falhe, o experimento seguirá projetado com pausas para registro individual.

A mudança de organização vai para o registro docente, porque a forma de acesso ao recurso pode alterar as interações observadas. Manter uma atividade pedagogicamente equivalente vale mais do que insistir em uma configuração técnica que a escola não consiga sustentar.

A mesma verificação confere se os estudantes compreendem os enunciados sem domínio prévio do jargão da área. Termos como característica e limiar, ou mesmo vizinho e centroide, entram sempre acompanhados de exemplos visíveis e explicações situadas. Caso alguma pergunta gere dúvida recorrente, a revisão do enunciado ou da mediação será tratada como ajuste necessário ao produto, afastando a ideia de falha do estudante para remover barreiras linguísticas alheias ao conteúdo.

\section{Alinhamento curricular preliminar}

Uma matriz explicitará o vínculo da sequência com a Computação na Educação Básica, associando cada habilidade curricular ao seu objetivo de aprendizagem, bem como à atividade proposta e às evidências observáveis. A versão preliminar apoia-se nas referências do Ensino Médio voltadas aos fundamentos da inteligência artificial com suas potencialidades e limites, articulando a análise em ciência de dados ao uso de modelos computacionais simples de simulação e previsão. Esse recorte dialoga com as habilidades EM13CO10 a EM13CO12 previstas no anexo da Resolução CNE/CEB nº 1/2022 \cite{brasil2022computacao}. Embora a norma deixe de prescrever esses cinco algoritmos nominalmente, eles oferecem situações didáticas concretas para examinar os limites de cada regra na transformação de dados em resultados. A matriz indicará ainda conexões com a Matemática, mantendo a Computação como objeto principal. O recorte definitivo das habilidades dependerá da série e da disciplina definidas, respeitando o currículo da rede estadual e as condições do campo.
